\section{Exploratory Data Analysis}
\label{sec:eda}

\subsection{Dataset construction}

The current candidate-presentation collection contains 22,355 sampled MuSiQue
questions. Each episode follows the same sequence: the Solver produces an
initial answer (A1), the Critic independently reviews that answer and supplies
a recommendation (C1), and the Solver produces a second answer (A2) after both
candidates are presented in randomized order. Generation is deterministic for
the recorded run. Solver activations are captured at decoder layers 4, 8, 17,
25, 29, 30, 31, 32, and 33 so that later analyses can compare where
feedback-related behavior is most decodable.

Questions are assigned to frozen train, validation, test, and intervention
partitions by question identifier. All records associated with one question
remain in the same partition. The training partition is used to fit SAEs and
probes, validation supports model and feature selection, test is reserved for
final predictive evaluation, and the intervention partition is reserved for
causal experiments. This question-level separation prevents related
interactions from appearing on both sides of an evaluation boundary.

\subsection{Behavior-label distribution}

The labeling policy compares A1, C1, and A2 using normalized text,
question-specific answer aliases, and conservative rules for refusals and
ambiguous responses. Each episode receives one of four mutually exclusive
labels:

\begin{itemize}
    \item \textbf{Adoption}: A1 and C1 differ, and A2 follows the Critic.
    \item \textbf{Retention}: A1 and C1 differ, and A2 retains A1.
    \item \textbf{Agreement}: A1 and C1 agree, and A2 preserves that answer.
    \item \textbf{Out}: the behavior cannot be compared reliably or the
    episode fails an inclusion rule.
\end{itemize}

\begin{table}[htbp]
    \centering
    \caption{Behavior labels in the candidate-presentation collection.}
    \label{tab:eda-label-counts}
    \begin{tabular}{lrr}
        \toprule
        Label & Episodes & Share \\
        \midrule
        Adoption  & 5,834 & 26.1\% \\
        Retention & 4,810 & 21.5\% \\
        Agreement & 8,451 & 37.8\% \\
        Out       & 3,260 & 14.6\% \\
        \midrule
        Total     & 22,355 & 100.0\% \\
        \bottomrule
    \end{tabular}
\end{table}

Adoption and retention together form the primary binary cohort of 10,644
episodes, or 47.6\% of the full collection. Within this cohort, adoption
accounts for 54.8\% and retention for 45.2\%. This is substantially more
balanced than an earlier free-form-feedback labeling approach, which was
dominated by adoption. Nevertheless, class-aware metrics such as balanced
accuracy, per-class precision and recall, F1 score, and area under the ROC
curve remain preferable to raw accuracy.

\begin{table}[htbp]
    \centering
    \caption{Behavior labels by frozen experimental partition.}
    \label{tab:eda-partitions}
    \begin{tabular}{lrrrr}
        \toprule
        Partition & Adoption & Retention & Agreement & Out \\
        \midrule
        Train        & 4,146 & 3,465 & 6,063 & 2,262 \\
        Validation   &   625 &   491 &   901 &   396 \\
        Test         &   529 &   416 &   755 &   292 \\
        Intervention &   534 &   438 &   732 &   287 \\
        Unassigned   &     0 &     0 &     0 &    23 \\
        \bottomrule
    \end{tabular}
\end{table}

The label proportions are broadly similar across the four frozen partitions,
which reduces the risk that an observed evaluation difference is driven only
by a major shift in the outcome distribution. The 23 unassigned questions are
quarantined as \emph{Out} and are not used for model fitting or evaluation.

\subsection{Data-quality observations}

Several response patterns require special handling. The collection includes
1,060 episodes in which both A1 and the Critic refuse or reject the premise;
these are generally agreements rather than adoption events. Other ambiguous
patterns include letter-only candidate selections, citation-only replies,
hyphenation differences, hybrid answers that combine both candidates, and
Critic responses that repeat A1 while nominally disagreeing. The final rules
retain explicit flags for these cases and exclude uncertain comparisons from
the primary binary target.

These observations constrain the interpretation of downstream probes. A probe
could learn answer length, refusal style, candidate position, or Critic
correctness instead of an internal tendency to accept feedback. Downstream
analysis should therefore include text-only and majority-class baselines,
report results by relevant behavioral flags, and verify that candidate order
is not predictive of the label. Probe-selected features should be interpreted
as correlates until targeted intervention demonstrates an effect beyond
random-feature and reconstruction controls.
